Personalized route recommendation has been studied with stronger data assumptions than this prototype currently supports. The closest conceptual baseline is Personalized Route Recommendation Based on User Habits for Vehicle Navigation, which models route sorting from historical navigation behavior and compares against minimum ETA, LightGBM, and DCN-v2-style methods using metrics such as mean inconsistency rate and AUC~\cite{userhabits2024}. That work motivates our use of historical route features and explicit baselines, but our setting differs: we focus on OSM walking routes and pseudo-history profiles derived from public trace signals rather than clean vehicle navigation logs.

Neural search methods frame personalized routing as an origin-destination route search problem. NASR integrates neural models with A* search for personalized route recommendation~\cite{nasr2019}. This is a stronger top-tier reference because it affects route generation itself, not only post-hoc ranking. Our current system is simpler: it generates a small diverse set of OSM candidate routes and reranks them with profile or prompt signals. A future version could compare against learned search or learned candidate-generation methods.

Trajectory reconstruction and popular-route mining are also relevant. Constructing Popular Routes from Uncertain Trajectories studies route construction from noisy or uncertain observations~\cite{chen2012popular}. Effective and Efficient Route Planning Using Historical Trajectories on Road Networks uses trajectory history to plan realistic routes efficiently~\cite{tian2023drpk}. These papers motivate two design requirements for our prototype: report reconstruction quality separately from ranking quality, and avoid treating noisy traces as ground truth without diagnostics.

Finally, preference-aware LLM planning work such as Personal Travel Solver shows how implicit preferences can be extracted and used in personalized travel planning~\cite{wang2025pts}. Our work is narrower: it ranks OSM walking-route candidates using structured route features and optional natural-language preferences. The connection is useful for evaluation language around preference satisfaction, but this paper should not claim full travel-planning capability.

\textbf{TODO:} Before submission, verify all citation metadata and copy exact metric names from each related paper.
